\chapter{Ejecución de los DSL}

En este capítulo se recopila la información de ejecución y validación para cada uno de los lenguajes de dominio específico desarrollados. Esta información se extrae directamente de los archivos \texttt{README.md} de cada proyecto DSL.

\section{DSL CIM Machine}

Este lenguaje permite la definición conceptual de máquinas de audio en el nivel \textit{CIM}. Define la identidad, puertos y configuraciones base de las máquinas.

\subsection{Comandos Principales}
\begin{itemize}
    \item \texttt{npm run langium:generate}: Genera el código del lenguaje a partir de la gramática.
    \item \texttt{npm run build}: Compila los servicios y validadores.
\end{itemize}

\subsection{Validación de Modelos JSON}
Se utiliza el script \texttt{validate.ts} para asegurar que los modelos cumplen tanto con la estructura \texttt{JSON} como con las reglas del lenguaje.

\begin{lstlisting}[language=bash]
# Validar un fichero individual
node validate.ts examples/valid/valid-1.json

# Validar todos los ficheros de un directorio
node validate.ts examples/valid
\end{lstlisting}

\section{DSL CIM Relations}

Este lenguaje define las conexiones conceptuales entre máquinas \textit{CIM}, permitiendo modelar el flujo de señales de audio y modificaciones a nivel abstracto.

\subsection{Comandos Principales}
\begin{itemize}
    \item \texttt{npm run langium:generate}: Genera el código del lenguaje.
    \item \texttt{npm run build}: Compila los servicios de validación.
\end{itemize}

\subsection{Validación de Modelos JSON}
Soporta validación interna y validación cruzada referenciando máquinas reales.

\begin{lstlisting}[language=bash]
# Validacion Interna
node validate.ts examples/valid

# Validacion Cruzada
node validate.ts examples/valid examples/machines
\end{lstlisting}

\section{DSL PIM Machine}

Este lenguaje permite definir máquinas de audio a nivel \textit{PIM}, incluyendo generadores, modificadores y efectos.

\subsection{Comandos Principales}
\begin{itemize}
    \item \texttt{npm run langium:generate}: Genera el código TypeScript a partir de la gramática.
    \item \texttt{npm run build}: Compila el proyecto.
\end{itemize}

\subsection{Validación de Modelos JSON}
\begin{lstlisting}[language=bash]
# Validacion individual o directorio
node validate.ts examples/valid

# Validacion Cruzada (PIM + Universo CIM)
node validate.ts examples/valid examples/machines
\end{lstlisting}

\section{DSL PIM Relations}

DSL para la orquestación de interconexiones técnicas entre máquinas \textit{PIM}. Permite establecer vínculos directos entre puertos de salida externos y parámetros o puertos de entrada.

\subsection{Validación}
Soporta dos modos: estructural (sintáctica) y cruzada (semántica).

\begin{lstlisting}[language=bash]
# Validacion Estructural
node validate.ts <fichero_o_directorio>

# Validacion Cruzada
node validate.ts <fichero_o_directorio> examples/machines
\end{lstlisting}
